% =====================================================================
%  FICHE 11 — THEME 4 — Corrigé DIFFICILE — Potentiels standard et prévision
% =====================================================================
\documentclass[11pt,a4paper]{article}
\usepackage[utf8]{inputenc}
\usepackage[T1]{fontenc}
\usepackage{lmodern}
\usepackage[french]{babel}
\usepackage{amsmath,amssymb}
\usepackage[version=4]{mhchem}
\usepackage{siunitx}
\sisetup{output-decimal-marker={,},per-mode=symbol}
\usepackage{xcolor}
\usepackage{array}
\usepackage{booktabs}
\usepackage{enumitem}
\usepackage[most]{tcolorbox}
\usepackage[a4paper,margin=2.0cm,headheight=26pt,headsep=10pt]{geometry}
\usepackage{fancyhdr}
\usepackage[hidelinks]{hyperref}

\definecolor{primary}{RGB}{146,95,160}
\definecolor{primarydark}{RGB}{92,52,104}
\definecolor{corrcol}{RGB}{0,135,90}
\newcommand{\NO}{\ensuremath{\mathrm{N.O.}}}
\newcommand{\Ered}{\ensuremath{E_{\mathrm{r}}^{\circ}}}

\newcounter{exo}
\newtcolorbox{corr}[1][]{enhanced,breakable,boxrule=0.9pt,arc=2pt,colback=corrcol!4,
  colframe=corrcol,fonttitle=\bfseries,coltitle=white,left=3mm,right=3mm,top=2.2mm,bottom=2.2mm,
  title={Corrigé~\refstepcounter{exo}\theexo\ifx\relax#1\relax\relax\else\ \textmd{--- \itshape #1}\fi}}

\pagestyle{fancy}\fancyhf{}
\fancyhead[L]{\small\textcolor{primarydark}{\textbf{Fiche 11 · Thème 4} — Corrigé}}
\fancyhead[R]{\small\textcolor{primarydark}{\textsc{Difficile}}}
\fancyfoot[C]{\small\thepage}
\renewcommand{\headrulewidth}{0.8pt}
\renewcommand{\headrule}{\hbox to\headwidth{\color{primary}\leaders\hrule height \headrulewidth\hfill}}

\begin{document}
\begin{center}
{\LARGE\bfseries\color{primarydark}Thème 4 — Corrigé Difficile}\\[-2pt]
{\color{corrcol}\rule{\textwidth}{1pt}}
\end{center}

\begin{corr}[une lame de zinc dans une solution mixte]
\begin{enumerate}[label=\textbf{\alph*)},leftmargin=2.2em,topsep=2pt,itemsep=2pt]
  \item Oxydants croissants : $\ce{Zn^{2+}}(-0{,}76)<\ce{Cu^{2+}}(+0{,}34)<\ce{Ag+}(+0{,}80)$.
        Réducteurs croissants (sens inverse) : $\ce{Ag}(+0{,}80)<\ce{Cu}(+0{,}34)<\ce{Zn}(-0{,}76)$,
        soit \ce{Zn} le plus fort réducteur.
  \item \ce{Cu^{2+}} (couple $+0{,}34$) oxyde les réducteurs des couples \emph{plus bas} : seul \ce{Zn}
        ($-0{,}76$) convient (\ce{Ag} est plus haut). $\ce{Cu^{2+} + Zn -> Cu + Zn^{2+}}$.
  \item Le zinc ($-0{,}76$) est plus bas que les deux couples : il réduit \textbf{les deux} ions.
        \[
        \ce{Zn + 2 Ag+ -> Zn^{2+} + 2 Ag} \qquad ; \qquad \ce{Zn + Cu^{2+} -> Zn^{2+} + Cu}.
        \]
  \item L'ion réduit en premier est le \textbf{meilleur oxydant}, soit \ce{Ag+} ($+0{,}80>+0{,}34$) :
        $\Delta E^{\circ}=0{,}80-(-0{,}76)=+1{,}56$ contre $+1{,}10$ pour \ce{Cu^{2+}}. L'argent se dépose avant le cuivre.
  \item Un bijou en \ce{Ag} appartient au couple $+0{,}80$, \emph{plus haut} que \ce{Cu^{2+}}/\ce{Cu}
        ($+0{,}34$) : $\Delta E^{\circ}<0$, donc \ce{Cu^{2+}} ne peut pas oxyder l'argent. Le bijou reste intact.
\end{enumerate}
\end{corr}

\begin{corr}[une médiamutation du manganèse]
\begin{enumerate}[label=\textbf{\alph*)},leftmargin=2.2em,topsep=2pt,itemsep=2pt]
  \item \ce{KMnO4} : $\NO(\ce{Mn})=+\mathrm{VII}$ ; \ce{MnSO4} (\ce{Mn^{2+}}) : $+\mathrm{II}$ ; \ce{MnO2} : $+\mathrm{IV}$.
  \item Le \ce{Mn} de \ce{KMnO4} passe de $+\mathrm{VII}$ à $+\mathrm{IV}$ : il est \textbf{réduit}.
        Le \ce{Mn} de \ce{MnSO4} passe de $+\mathrm{II}$ à $+\mathrm{IV}$ : il est \textbf{oxydé}.
        Deux degrés qui convergent vers un intermédiaire : c'est une \textbf{médiamutation}.
  \item \ce{KMnO4} est \textbf{réduit}, donc c'est l'\textbf{oxydant}. La justification de l'étudiant est
        \emph{fausse} : un oxydant \textbf{capte} des électrons, il ne les \emph{cède} pas.
        (La conclusion « oxydant » est juste, mais pour la mauvaise raison — annales 2020-Q1, prop. B.)
  \item Pour que \ce{MnO4^-} oxyde \ce{Mn^{2+}}, il faut :
        \[
        \Ered(\ce{MnO4^-}/\ce{MnO2}) > \Ered(\ce{MnO2}/\ce{Mn^{2+}}) \quad(\text{soit } \Delta E^{\circ}>0),
        \]
        ce qui est le cas ici : c'est la proposition \textbf{D} des annales. \emph{Équation équilibrée :}
        $\ce{2 KMnO4 + 3 MnSO4 + 2 H2O -> 5 MnO2 + K2SO4 + 2 H2SO4}$.
\end{enumerate}
\end{corr}

\begin{corr}[trois réactions : lesquelles ont lieu ?]
On calcule $\Delta E^{\circ}=\Ered(\text{oxydant}) - \Ered(\text{réducteur})$ ; la réaction a lieu si $\Delta E^{\circ}>0$.
\begin{enumerate}[label=\textbf{\alph*)},leftmargin=2.2em,topsep=2pt,itemsep=2pt]
  \item Oxydant \ce{S4O6^{2-}} ($+0{,}08$), réducteur \ce{I^-} (couple \ce{I2}/\ce{I^-}, $+0{,}536$) :
        $\Delta E^{\circ}=0{,}08-0{,}536=-0{,}456<0$ $\Rightarrow$ \textbf{n'a pas lieu}. Le sens spontané est l'inverse :
        \ce{I2} oxyde \ce{S2O3^{2-}} (c'est la base de l'iodométrie).
  \item Oxydant \ce{Cr2O7^{2-}} ($+1{,}23$), réducteur \ce{Sn^{2+}} (couple \ce{Sn^{4+}}/\ce{Sn^{2+}}, $+0{,}151$) :
        $\Delta E^{\circ}=1{,}23-0{,}151=+1{,}08>0$ $\Rightarrow$ \textbf{a lieu}.
  \item Oxydant \ce{Br2} ($+1{,}06$), réducteur \ce{Cl^-} (couple \ce{Cl2}/\ce{Cl^-}, $+1{,}39$) :
        $\Delta E^{\circ}=1{,}06-1{,}39=-0{,}33<0$ $\Rightarrow$ \textbf{n'a pas lieu}. C'est \ce{Cl2} qui oxyde
        \ce{Br^-} (sens inverse spontané).
\end{enumerate}
\end{corr}

\end{document}
